\section{The Container Confinement Problem}
\label{sec:container-confinment-problem}


Container confinement can be defined as the set of mechanisms and strategies used to isolate a containerized application from the host system and other containers. Effective confinement is critical for maintaining system security and stability, as it ensures that containers operate within defined boundaries by restricting their access to system resources, processes, networks, inter-process communication (IPC), filesystems, devices, and kernel interfaces. 

A challenge with container confinement is that the level of confinement of a container is not fixed; it can vary based on the specific requirements or intentions of the user. In other words, the degree to which a container is isolated from the host system or other containers is customizable to meet the user’s needs. At one end of the spectrum, one can lock a container down so tightly that it is almost entirely isolated from the host system, which minimizes potential attack surfaces but can also render the container incapable of performing key tasks. On the other hand, allowing a container too much freedom undermines the security benefits of containerization by exposing more privileges than necessary. The optimal configuration, however, is one that strikes a balance between maximum confinement (restricting access to system resources as much as possible) and function flexibility (enabling the container to perform its intended tasks). Achieving this balance is the core problem and goal of container confinement.

A standard methodology for assessing container confinement is crucial due to the complexity of modern systems and the intricate vulnerabilities they present. A formal approach involving rigorously evaluating and analyzing the security of the container system might seem ideal, however, it is impractical because container environments are highly complex, with interactions between kernels, runtimes, orchestration tools, and external dependencies that are difficult to model comprehensively. Moreover, many vulnerabilities stem from subtle details that a formal model might overlook. Instead, an empirical approach offers a practical alternative. Though it will never be exhaustive, the goal is not comprehensive coverage but rather establishing a baseline set of tests that demonstrate effective container confinement. This approach is achievable and can evolve over time as new vulnerabilities are identified, allowing the methodology to remain current. Starting with tests for isolation, privilege escalation prevention, and resource limitation enforcement provides a solid foundation. By focusing on demonstrating security effectiveness in critical areas, this methodology can enhance confidence in container confinement while enabling continuous improvement and adaptation to the evolving threat landscape.